\subsection{The Galois group of the algebraic closure of a finite field}

Let \(S \neq \{1\}\) be a set of integers \(\geqslant 1\) stable under multiplication; we shall order it by the relation « \(m\) divides \(n\) ». When \(m\) divides \(n\) , we have \(mZ \supset nZ\) , hence there is a canonical homomorphism \(\pi_{m,n}\) of the ring \(Z/nZ\) onto the ring \(Z/mZ\) . We denote by \(A(S)\) the inverse limit of the inverse system of rings \((Z/mZ, \pi_{m,n})\) indexed by \(S\) . Each finite set \(Z/mZ\) is equipped with the discrete topology and \(A(S)\) with the topology induced by that of the product \(\prod_{n \in S} (Z/nZ)\) .

Then A(S) is a compact topological ring (Gen.~Top. I, p.~64, Prop. 8). We see at once that the unique ring homomorphism \(\varphi\) of \(\mathbf{Z}\) into A(S) is injective with a dense image; we shall identify \(\mathbf{Z}\) with its image under \(\varphi\) in A(S). For the topology induced on \(\mathbf{Z}\) by that of A(S), the sets \(m\mathbf{Z}\) (for \(m \in S\) ) form a base of neighbourhoods of 0.

When \(S = N^{*}\) , \(A(S)\) is written \(\hat{Z}\) . When S consists of the powers of a prime number l, \(A(S)\) is written \(Z_{l}\) and is called « ring of l-adic integers ». We thus have

\[
\hat {\mathbf {Z}} = \varprojlim_ {m \geqslant 1} \mathbf {Z} / m \mathbf {Z}, \quad \mathbf {Z} _ {l} = \varprojlim_ {n \geqslant 0} \mathbf {Z} / l ^ {n} \mathbf {Z}.
\]

When S and T are two sets of integers stable under multiplication such that \(S \supset T\) , we have a natural projection \(\mathbf{A}(\mathbf{S}) \to \mathbf{A}(\mathbf{T})\) which is a continuous homomorphism of topological rings. In particular, for every prime number l we have a continuous homomorphism \(\hat{Z} \to Z_{l}\) . From it we obtain a continuous homomorphism

\[
\hat {\mathbf {Z}} \rightarrow \prod_ {l} \mathbf {Z} _ {l}
\]

(product extended over all prime numbers) ; this is an isomorphism of topological rings, as follows from the passage to the inverse limit in I, p.~112, Prop. 11.

Let K be a finite field of q elements and \(\bar{K}\) an algebraic closure of K. For every integer \(m \geqslant 1\) we denote by \(K_{m}\) the unique subfield of \(\bar{K}\) which is of degree m over K (Prop. 3). We have \(\bar{K} = \bigcup_{m \geqslant 1} K_{m}\) . Further we denote by \(\sigma_{q}\) the automorphism

\(x \mapsto x^{q}\) of the perfect field \(\bar{K}\) ; it is called the Frobenius automorphism of \(\bar{K}\) (relative to K).

PROPOSITION 5. --- There exists a unique isomorphism of topological groups \(\pi_{\mathbf{K}}: \hat{\mathbf{Z}} \to \operatorname{Gal}(\bar{\mathbf{K}} / \mathbf{K})\) such that \(\pi_{\mathbf{K}}(1) = \sigma_q\) .

Let \(\Gamma\) be the subgroup of \(\operatorname{Gal}(\bar{\mathbf{K}} / \mathbf{K})\) generated by \(\sigma_q\) . For every integer \(m > 0\) we have \(\sigma_q^m (x) = x^{q^m}\) for all \(x\in \bar{\mathbf{K}}\) , hence the set of fixed points of \(\sigma_q^m\) is equal to \(\mathbf{K}_m\) . Since \(\mathbf{K}_m\neq \bar{\mathbf{K}}\) , we have \(\sigma_q^m\neq 1\) . Hence there is an isomorphism \(\pi_0\) of \(\mathbf{Z}\) onto \(\Gamma\) which maps 1 to \(\sigma_q\) .

The field of invariants of \(\Gamma\) consists of the \(x \in \bar{K}\) such that \(x^{q} = x\) , hence is equal to K. Therefore (V, p.~67, Lemma 2) the group \(\Gamma\) is dense in \(\operatorname{Gal}(\bar{\mathbf{K}}/\mathbf{K})\) . Since every subextension of \(\bar{K}\) of finite degree over K is one of the fields \(K_{m}\) , the subgroups \(\operatorname{Gal}(\bar{\mathbf{K}}/\mathbf{K}_{m})\) form a fundamental system of neighbourhoods of 1 in \(\operatorname{Gal}(\bar{\mathbf{K}}/\mathbf{K})\) . It is clear that \(\Gamma \cap \operatorname{Gal}(\bar{\mathbf{K}}/\mathbf{K}_{m})\) is the cyclic group generated by \(\sigma_{q}^{m}\) , hence equal to \(\pi_{0}(m\mathbf{Z})\) .

From the above remarks about the topology of \(\hat{Z}\) the isomorphism \(\pi_{0}:Z\to\Gamma\) extends in a unique fashion to an isomorphism of topological groups \(\pi_{K}:\hat{Z}\to\operatorname{Gal}(\bar{\mathbf{K}}/\mathbf{K})\) .

Let \(m \geqslant 1\) be an integer; it is clear that the Frobenius automorphism of \(\bar{\mathbf{K}}\) relative to \(\mathbf{K}_m\) is \(\sigma_q^m\) . Hence we obtain the relation

\[
\pi_ {\mathrm{K} _ {m}} (a) ^ {*} = \pi_ {\mathrm{K}} (m a) \quad \text { for } \quad a \in \hat {\mathbf {Z}}.\tag{5}
\]

